\phantomsection\label{fig:vol01:western-zhou:overview-map}
\section{空间总览：王畿、东方与边疆}
\begin{center}
\begin{tikzpicture}[x=.82cm,y=.82cm,font=\small]
  \fill[paper] (0,0) rectangle (12,7);
  \fill[blue!8] (10.7,0) rectangle (12,7);
  \node[rotate=90,text=timeline] at (11.35,3.5) {东海方向};

  % Rivers are orientation tools rather than surveyed hydrography.
  \draw[blue!60,line width=1pt] (0.2,4.9) .. controls (3,4.6) and (4.7,5.3) .. (8.4,4.7);
  \node[text=blue!70!black,above] at (4.3,5.0) {黄河};
  \draw[blue!55,line width=.8pt] (4.2,.7) .. controls (5.2,2.1) and (6.5,2.2) .. (8.6,1.2);
  \node[text=blue!70!black,below] at (6.4,1.55) {汉水流域};

  \fill[dynasty!25] (2.1,3.5) -- (4.2,3.1) -- (4.6,4.3) -- (3.2,5.0) -- cycle;
  \node[align=center] at (3.4,4.1) {周王畿\\镐京};
  \fill[timeline!20] (5.2,3.3) -- (6.8,3.5) -- (7.0,4.6) -- (5.8,4.9) -- cycle;
  \node[align=center] at (6.1,4.1) {成周\\洛邑};
  \fill[orange!20] (7.4,4.6) -- (9.0,4.8) -- (9.4,5.9) -- (7.8,6.0) -- cycle;
  \node at (8.5,5.35) {燕};
  \fill[green!18] (8.2,3.0) -- (9.7,3.1) -- (10.0,4.2) -- (8.8,4.5) -- cycle;
  \node[align=center] at (9.1,3.75) {齐鲁\\东方封国};
  \fill[orange!25] (4.5,1.7) -- (6.2,1.4) -- (6.4,2.7) -- (5.1,3.0) -- cycle;
  \node[align=center] at (5.4,2.2) {申及\\南阳方向};
  \fill[gray!18] (.6,2.3) -- (2.0,2.6) -- (2.1,3.7) -- (.9,4.0) -- cycle;
  \node[align=center] at (1.25,3.25) {西部\\戎狄诸部};
  \fill[purple!15] (6.9,.4) -- (9.5,.6) -- (9.9,2.6) -- (8.1,2.9) -- cycle;
  \node[align=center] at (8.4,1.6) {南方政治\\共同体};

  \draw[-{Stealth[length=2mm]},thick,color=dynasty] (3.9,4.1) -- (5.5,4.1);
  \node[above,text=dynasty] at (4.7,4.15) {东向经营};
  \draw[-{Stealth[length=2mm]},thick,color=timeline] (3.6,3.55) -- (6.8,1.7);
  \node[below,text=timeline] at (5.2,2.9) {南征方向};
\end{tikzpicture}

\smallskip
\small\textit{图  西周空间关系示意：突出王畿、成周、东方封国与南方/西部边疆的相对位置；并非精确疆界图。参考传世文献、金文研究与历史地理通说。}
\end{center}
